\chapter{Protocolo}

Este capítulo explica lo implementado del protocolo: las peticiones del comprador, los turnos con los agentes de las partes y el cierre con \texttt{SESSION\_CLOSE}. El mensaje \texttt{AGREEMENT\_HELD} está pendiente de implementar.

\section{Peticiones del comprador}

El coordinador del comprador es cliente A2A del orquestador (D-87). Le envía sus peticiones en el Message de \texttt{SendMessage}, y el orquestador responde siempre con una Task. Lo único implementado es \texttt{OPEN\_ROUND}, en su primera ronda.

\subsection{Autenticación}

Cada petición lleva \texttt{Authorization: Bearer <token>} (D-128). La AgentCard del orquestador es pública y declara ese esquema. Sin un token válido, el orquestador responde 401 con \texttt{WWW-Authenticate}. Es el único caso en el que no queda rastro en el registro de peticiones: así, nadie sin credencial puede llenar las tablas.

\subsection{Forma de la petición}

El Message lleva la cabecera \texttt{A2A-Extensions}, la URI de la extensión en \texttt{extensions}, el sobre en \texttt{metadata} y una sola Part \texttt{data} (D-19, D-103). Este es el contenido, sin la configuración de cada sesión:

\begin{verbatim}
metadata["urn:tfm:negociacion:v1"]:
  { "type": "OPEN_ROUND", "sender": "comprador-1" }

Part data, con mediaType
application/vnd.tfm.negociacion.peticion+json:
  { "sessions": [ { "itemId": "partida-1", "config": { ... } } ] }
\end{verbatim}

\subsection{Resultado}

Si la petición es válida, la Task empieza en \texttt{TASK\_STATE\_WORKING} y se completa cuando todas las sesiones han salido de \texttt{CREATED} (D-94). Con \texttt{returnImmediately = true}, el coordinador recibe la Task enseguida y la consulta con \texttt{GetTask}; sin él, la respuesta espera al resultado. El Artifact dice qué sesiones han empezado:

\begin{verbatim}
{ "requestId": "...", "round": 1,
  "sessions": [
    { "itemId": "partida-1", "sessionId": "...",
      "started": true },
    { "itemId": "partida-2", "sessionId": "...",
      "started": false, "failedPartyId": "vendedor-2" } ] }
\end{verbatim}

\subsection{Rechazos}

Una petición con fallos no crea nada: la Task termina en \texttt{TASK\_STATE\_REJECTED} y su mensaje de estado lleva todos los fallos, ordenados por ruta (D-89, D-101). Los índices se ordenan como números, así que \texttt{sessions.2} va antes que \texttt{sessions.10}.

\begin{verbatim}
{ "rejections": [
    { "code": "INVALID_CONFIG",
      "field": "body.sessions.0.config.maxTurns" } ] }
\end{verbatim}

\begin{tabular}{p{4.2cm}p{9.3cm}}
  \textbf{Código} & \textbf{Cuándo} \\
  \texttt{INVALID\_MESSAGE} & El Message no tiene forma de petición (D-103). \\
  \texttt{MISSING\_FIELD} \newline \texttt{WRONG\_TYPE} \newline \texttt{UNKNOWN\_VALUE} \newline \texttt{UNKNOWN\_FIELD} & Fallos de estructura del sobre o del cuerpo. \\
  \texttt{INVALID\_CONFIG} & Una configuración incumple sus restricciones, hay un \texttt{itemId} repetido, el \texttt{BUYER} cambia entre sesiones o un origen no está permitido (D-102, D-129). \\
  \texttt{NOT\_AUTHORIZED} & El \texttt{sender} no es la identidad autenticada o no es el \texttt{BUYER} de las sesiones (D-102, D-128). \\
  \texttt{UNKNOWN\_REQUEST} & La petición lleva \texttt{requestId}. Es un límite de esta versión, que no implementa la ronda 2. \\
\end{tabular}

El mensaje de estado lleva un sobre con \texttt{sender = ORCHESTRATOR}. Si no se pudo leer un \texttt{type} de petición válido, va sin sobre, porque no hay un valor correcto que poner (D-106).

\section{Turnos con los agentes}

En cada turno, el orquestador es cliente A2A del agente de la parte con turno (D-05, D-06). Le pide un acto con un \texttt{TURN\_REQUEST} y el agente lo devuelve en el Artifact de una Task. El flujo completo está en el capítulo «Flujo de una sesión».

\subsection{\texttt{TURN\_REQUEST}}

Va en el Message de \texttt{SendMessage}, con \texttt{role = ROLE\_USER}, la URI de la extensión en \texttt{extensions} y en la cabecera \texttt{A2A-Extensions}, y \texttt{returnImmediately = true} (D-58). Lleva todo lo que el agente necesita, porque el agente no guarda nada de la sesión (D-46). Este es el del turno 2:

\begin{verbatim}
metadata["urn:tfm:negociacion:v1"]:
  { "sessionId": "...", "turn": 2,
    "type": "TURN_REQUEST", "sender": "ORCHESTRATOR" }

Part data, con mediaType
application/vnd.tfm.negociacion.control+json:
  { "role": "SELLER",
    "fixedTerms": { "productId": "p1", "unit": "kg",
                    "currency": "EUR" },
    "allowedActs": ["COUNTEROFFER", "ACCEPT", "WITHDRAW"],
    "turnsLeft": 9, "attempt": 1,
    "lastOffer": { "productId": "p1", "unit": "kg",
                   "currency": "EUR", "quantity": 10,
                   "unitPrice": 12.50, "deliveryDays": 30,
                   "paymentDays": 60 } }
\end{verbatim}

En el turno 1, \texttt{allowedActs} es \texttt{[PROPOSE, WITHDRAW]} y no hay \texttt{lastOffer} (D-47). \texttt{turnsLeft} cuenta el turno actual: \texttt{maxTurns − turn + 1}. El \texttt{messageId} es el guardado para el intento (D-66). El \texttt{contextId} solo va desde el segundo turno de cada parte: en el primero, el orquestador guarda el que devuelve el agente (D-20).

\subsection{El acto}

El agente responde con una Task que, al completarse, lleva un solo Artifact con una sola Part \texttt{data} de \texttt{mediaType} \texttt{application/vnd.tfm.negociacion.acto+json} (D-19, D-23). El sobre del Artifact devuelve el \texttt{sessionId} y el \texttt{turn} recibidos, con el tipo del acto y su \texttt{partyId} como \texttt{sender} (D-28, D-30):

\begin{verbatim}
metadata["urn:tfm:negociacion:v1"]:
  { "sessionId": "...", "turn": 2,
    "type": "COUNTEROFFER", "sender": "vendedor-1" }

Part data:
  { "offer": { ..., "unitPrice": 11.75, ... } }
\end{verbatim}

La validación es siempre estricta en esta versión (D-11): cualquier fallo de R-01 a R-17 se registra sin corregir y aborta la sesión.

\subsection{Números exactos}

El precio es un decimal exacto (D-37). Los tipos de \texttt{a2a-sdk} guardan los números como coma flotante y perderían cifras, así que el orquestador usa su propio cliente JSON-RPC (D-130). Lee los decimales como \texttt{Decimal} y los escribe con todas sus cifras, tanto hacia el agente como en PostgreSQL.

Por la misma razón, un agente hecho con el SDK envía el entero 10 como 10.0. En los campos enteros (\texttt{turn}, \texttt{quantity}, \texttt{deliveryDays} y \texttt{paymentDays}), un número con parte decimal nula vale como entero, sin coerción, y se reenvía como entero (D-132). Un entero de 4300 cifras o más es \texttt{OUT\_OF\_RANGE}: no se llega a convertir, porque eso agotaría la memoria del orquestador (D-133).

\subsection{Cabeceras y métodos}

El cliente llama a \texttt{SendMessage}, \texttt{GetTask} y \texttt{CancelTask} por JSON-RPC, con las cabeceras \texttt{A2A-Extensions: urn:tfm:negociacion:v1} y \texttt{A2A-Version: 1.0}. Sin la segunda, un servidor hecho con el SDK supone la versión 0.3 del protocolo. No sigue redirecciones: la URL que llama es la de la AgentCard, ya comprobada contra los orígenes permitidos (D-129, D-131).

\section{\texttt{SESSION\_CLOSE}}

Cuando la sesión entra en \texttt{CLOSING}, el orquestador comunica el resultado a las dos partes (CU-02). Se envía como el \texttt{TURN\_REQUEST}, en el Message de \texttt{SendMessage} con \texttt{returnImmediately = true}, pero con un \texttt{messageId} nuevo que no se guarda (D-64). El \texttt{contextId} solo va si la parte lo tiene: si la sesión cierra en el turno 1, la parte que no abre nunca lo recibió. El \texttt{turn} del sobre es el del turno que provocó el cierre.

\begin{verbatim}
metadata["urn:tfm:negociacion:v1"]:
  { "sessionId": "...", "turn": 3,
    "type": "SESSION_CLOSE", "sender": "ORCHESTRATOR" }

Part data:
  { "outcome": "WITHDRAWN", "withdrawnBy": "comprador-1",
    "lastOffers": [
      { "partyId": "comprador-1", "offer": { ... } },
      { "partyId": "vendedor-1", "offer": { ... } } ] }
\end{verbatim}

El cuerpo lleva \texttt{outcome} y solo los campos que le corresponden (AC-02-03):

\begin{tabular}{p{3.5cm}p{10cm}}
  \textbf{\texttt{outcome}} & \textbf{Campos} \\
  \texttt{AGREEMENT} & \texttt{agreedOffer}, sin \texttt{lastOffers}. \\
  \texttt{WITHDRAWN} & \texttt{withdrawnBy}. El \texttt{reason} del \texttt{WITHDRAW} no viaja: queda en el registro de actos aceptados. \\
  \texttt{NO\_AGREEMENT} & Ninguno más. \\
  \texttt{ABORTED} & \texttt{abortReason = RETRIES\_EXHAUSTED} y \texttt{lastFailure}. \\
  \texttt{RELEASED} & \texttt{releaseReason}. \\
\end{tabular}

Salvo en \texttt{AGREEMENT}, el cuerpo lleva \texttt{lastOffers} si alguna parte llegó a ofertar (D-82). \texttt{AGREEMENT} y \texttt{RELEASED} solo aparecen cuando la sesión sale de \texttt{HELD}, que aún no está implementado.

La respuesta no se valida (D-49): una Task en cualquier estado o un Message marcan el envío como enviado. No se sondea ni se envía \texttt{CancelTask}. Un error de red, la falta de respuesta en \texttt{turnTimeoutSeconds} o un error A2A lo marcan como fallido, con \texttt{AGENT\_UNREACHABLE}, \texttt{TURN\_TIMEOUT} o \texttt{AGENT\_ERROR}, sin reintento (D-64, D-71). Una respuesta que no es Task ni Message no está en la especificación: mientras no se decida, se anota como \texttt{AGENT\_ERROR}.
